\lecture{9}
\title{Inner Products, Orthonormality}

\begin{document}

Recall from vector algebra that the dot product tells us a lot about the lengths of vectors and angles between vectors.
\begin{quote}
    \textbf{Fact (Vector Algebra).} For any $\vec{x}, \vec{y} \in \mathbb{R}^n$, with angle $\theta$ between them, we have $\vec{x} \cdot \vec{y} = |\vec{x}| \cdot |\vec{y}| \cdot \cos \theta$.
\end{quote}
Let's generalize this in 18.701 terms.

\subsection{Inner Products in $\mathbb{R}^n$}

\textbf{Definition (Dot / Inner Product).} For vectors $x, y \in \mathbb{R}^n$, we say their \textbf{dot / inner product} is $\langle x, y \rangle := x^{\top}y \in \mathbb{R}$.

\textbf{Definition (Length / Norm / Distance).} For vectors $x \in \mathbb{R}^n$, the \textbf{length / norm} of $x$ is $|x| := \sqrt{\langle x, x \rangle}$.  Furthermore, the \textbf{distance} between $x, y \in \mathbb{R}^n$ is said to be $|x - y|$.

Note that, for now, we are working only over $\mathbb{R}^n$.  We don't want to say $\begin{psmallmatrix} 1 \\ i \end{psmallmatrix}$ has length $0$, for example---and some fields don't even have a well-defined square root!

\begin{remark}
    There's also something called an \textbf{outer product} defined by $x \otimes y := x y^{\top} \in \mathbb{R}^{n \times n}$.
\end{remark}

\begin{theorem}{Cauchy-Schwarz}
    For all $x, y \in \mathbb{R}^n$, we have $|\langle x, y \rangle | \leq |x| \cdot |y|$.
\end{theorem}

\begin{proof}
    The identity $|x + \lambda y|^2 \geq 0$ holds for all $\lambda \in \mathbb{R}$.  Expanding this, we have:
    \[ 0 \leq |x + \lambda y|^2 = \langle x + \lambda y, x + \lambda y \rangle = |x|^2 + 2\langle x, y \rangle \lambda + |y|^2 \lambda^2. \]
    The RHS of the above is a quadratic in $\lambda$ that is nonnegative everywhere, so its discriminant is nonpositive:
    \[ 0 \geq \text{discriminant} = \left ( 2 \langle x, y \rangle \right ) ^2 - 4 \left ( |x|^2 \right ) \left ( |y|^2 \right ) = 4 ( \langle x, y\rangle^2 - |x|^2 |y|^2). \]
    The above is just the statement of Cauchy-Schwarz. ~ $\blacksquare$
\end{proof}

\textbf{Definition (Angle).} The angle between $x, y \in \mathbb{R}^n$ is $\cos^{-1}\left(\frac{\langle x, y \rangle}{|x| \cdot |y|}\right)$ (which is well-defined by Cauchy-Schwarz).

\subsection{Orthonormal Bases and Matrices}

\textbf{Definition (Orthogonal).} Vectors $x$ and $y$ are \textbf{orthogonal} if $\langle x, y \rangle = 0$, meaning the angle between them is $90^{\circ}$.

\textbf{Definition (Orthogonal / Orthonormal Bases).} A basis of $\mathbb{R}^n$ is \textbf{orthogonal} if all vectors in the basis are pairwise orthogonal.  A basis of $\mathbb{R}^n$ is \textbf{orthonormal} if it is orthogonal and all vectors in the basis have unit length.

\begin{theorem}{Orthogonal $\implies$ Independent}
    Any set of pairwise orthogonal vectors $\{a_1, \dots, a_k \} \subseteq \mathbb{R}^n$ must be linearly independent.  In particular, $n$ pairwise orthogonal vectors of $\mathbb{R}^n$ must form a basis.
\end{theorem}

\begin{proof}
    Given any linear independence of $\{a_1, \dots, a_k\}$, we may write:
    \[ c_1a_1 + \dots + c_k a_k  = 0 ~ \implies ~ \langle a_i, c_1 a_1 + \dots + c_k a_k \rangle = 0 \text{ for all } i = 1, 2, \dots, n. \]
    But $\langle a_i, c_1 a_1 + \dots + c_k a_k \rangle = c_i |a_i|^2$, so $c_i = 0$ for all $i$. ~ $\blacksquare$
\end{proof}

\textbf{Definition (Orthogonal Matrices).} A matrix $A \in \mathbb{R}^{n \times n}$ is \textbf{orthogonal} if $\langle Ax, Ay \rangle = \langle x, y \rangle$ for all $x, y \in \mathbb{R}^n$.

In other words, an orthogonal matrix is a matrix that preserves the dot product.

\begin{theorem}{Properties of Orthogonal Matrices}
    For any $A \in \mathbb{R}^{n \times n}$, the following are equivalent:
    \begin{enumerate}
        \item The matrix $A$ is orthogonal.
        \item For all $x \in \mathbb{R}^n$, we have $|Ax| = |x|$.
        \item We have $A^{\top}A = I_n$.
        \item The columns of $A$ form an orthonormal basis.
        \item The rows of $A$ form an orthonormal basis.
    \end{enumerate}
\end{theorem}

\begin{proof}
    We just have to show a bunch of implications.
    \begin{itemize}
        \item $(\#1) \implies (\#2)$.  If $A$ is orthogonal, then $|Ax| = \sqrt{\langle Ax, Ax \rangle} = \sqrt{\langle x, x \rangle} = |x|$.
        \item $(\#2) \implies (\#1)$.  This follows from the polarization identity:
        \[ \langle x, y \rangle = \dfrac{|x + y|^2 - |x|^2 - |y|^2}{2}. \]
        Since $A$ preserves lengths, it preserves the RHS, so it must also preserve the LHS.
        \item $(\#3) \implies (\#1)$.  If $A^{\top}A = I_n$, then $\langle Ax, Ay \rangle = (Ax)^{\top}(Ay) = (x^{\top}A^{\top})(Ay) = x^{\top}(A^{\top}A)y = x^{\top}y = \langle x, y \rangle$.
        \item $(\#1) \implies (\#3)$.  Take a pair of basis vectors $(e_i, e_j)$, so $\langle Ae_i, Ae_j\rangle = \langle e_i, e_j \rangle$.  The LHS and RHS simplify as:
        \[ \text{LHS} = \langle Ae_i, Ae_j \rangle = e_i^{\top}(A^{\top}A)e_j = (A^{\top}A)_{ij} ~~~~ \text{ and } ~~~~ \text{RHS} = \langle e_i, e_j \rangle = \begin{cases} 1 & i = j \\ 0 & i \neq j \end{cases}. \]
        Thus, the equation $\langle Ae_i, Ae_j\rangle = \langle e_i, e_j \rangle$ gives us $(A^{\top}A)_{ij} = 1$ if and only if $i = j$, meaning $(A^{\top}A) = I_n$.
        \item $(\#4) \iff (\#2)$.  Suppose the columns of $A$ form vectors $\{a_1, \dots, a_n\} \subseteq \mathbb{R}^n$.  Observe that:
        \[ A^{\top}A = \begin{bsmallmatrix} - & a_1 & - \\ \vdots & \vdots & \vdots \\ - & a_n & - \end{bsmallmatrix} \begin{bsmallmatrix} \mid & \cdots & \mid \\ a_1 & \cdots & a_n \\ \mid & \cdots & \mid \end{bsmallmatrix} = \begin{bmatrix} \langle a_i, a_j \rangle \end{bmatrix}_{i, j}. \]
        Thus, if $A^{\top}A = I_n$, then $\langle a_i, a_j \rangle = 1$ if $i = j$, and $\langle a_i, a_j \rangle = 0$ otherwise, which implies $\{a_i\}$ is an orthonormal basis.  And it goes the other way, too; if $\{a_i\}$ is orthonormal, then $A^{\top}A = I_n$.
        \item $(\# 5) \iff (\#2)$. If $A^{\top} A = I_n$, then $A^{\top}$ and $A$ are inverses.  So we also get $A A^{\top} = I_n$.  Now repeating the argument in $(\#2) \implies (\#4)$ with the equation $AA^{\top} = I_n$ implies the rows of $A$ are orthonormal.
    \end{itemize}
    And that's good enough. ~ $\blacksquare$
\end{proof}

\subsection{Orthonormal Groups}

\textbf{Definition (Orthogonal Group).} The \textbf{orthogonal group} is $O_n(\mathbb{R}) := \{A \in \mathbb{R}^{n \times n} \mid A^{\top} A = I_n \}$.

\textbf{Definition (Special Orthogonal Group).} The \textbf{special orthogonal group} is $SO_n(\mathbb{R}) := \{M \in O_n(\mathbb{R}) \mid \det M = 1\}$. 

Note that for all $M \in O_n(\mathbb{R})$, the condition $M^{\top}M = I_n$ implies $\det M = \pm 1$.  So in fact $[O_n(\mathbb{R}): SO_n(\mathbb{R})] = 2$.

For small $n$, what do the orthogonal groups look like?

\textbf{Example.} Evidently $O_1(\mathbb{R}) \cong \{ \begin{bsmallmatrix} 1 \end{bsmallmatrix} , \begin{bsmallmatrix} -1 \end{bsmallmatrix} \}$, and $SO_1(\mathbb{R}) \cong \{\begin{bsmallmatrix} 1 \end{bsmallmatrix} \}$.

\textbf{Example.} By inspection, for $n = 2$, we can describe the elements of $O_2$ and $SO_2$ as rotations:
\[ O_2(\mathbb{R}) = \left \{ \begin{bsmallmatrix} \cos \theta & - \sin \theta \\ \sin \theta & \cos \theta \end{bsmallmatrix}, \begin{bsmallmatrix} \cos \theta & \sin \theta \\ \sin \theta & - \cos \theta \end{bsmallmatrix} : \theta \in [0, 2\pi) \right \} \text{ and } SO_2(\mathbb{R}) = \left \{ \begin{bsmallmatrix} \cos \theta & -\sin \theta \\ \sin \theta & \cos \theta \end{bsmallmatrix}: \theta \in [0, 2\pi) \right \}. \]
Notice that $SO_2(\mathbb{R})$ is just the group of rotation matrices.  Meanwhile, $O_2(\mathbb{R})$ includes rotations about the origin \emph{or} reflections about lines through the origin.

\begin{quote}
    \textcolor{red!75}{\emph{Important!}} \ \textcolor{black!60}{Composing reflections and rotations always yields either a pure reflection or pure rotation!}
\end{quote}

\textbf{Example.} For any odd $n$ and any $A \in SO_n(\mathbb{R})$, we have:
\begin{align*}
\det (A - I_n) = \ &  \det (A - AA^{\top}) \\
= \ &  \det (A) \det (I_n - A^{\top}) ~~~ \text{ by factoring} \\
= \ & \det (I_n - A^{\top}) ~~~ \text{ since } \det(A) = 1\\
= \ & \det ((I_n - A)^{\top}) \\
= \ & \det (I_n - A) ~~~ \text{ since } \det(M^{\top}) = \det(M) \text{ for any } M \\
= \ & (-1)^n \cdot \det(A - I_n)
\end{align*}
But when $n$ is odd, the above means $\det (A - I_n) = 0$, meaning $A$ has an eigenvector $v$ with eigenvalue $1$.  In the particular case of $n = 3$, if we rewrite $A$ using an orthonormal basis containing $v = e_1$, we must have:
\[ Ae_1 = e_1 ~  \implies ~ A = \begin{bsmallmatrix} 1 & 0 & 0 \\ 0 & \heartsuit & \heartsuit \\ 0 & \heartsuit & \heartsuit \end{bsmallmatrix}. \]
Note that the top entries of $A$ must be zero because the columns of $A$ must be orthonormal.  And the submatrix of $\heartsuit$s is in $SO_2(\mathbb{R})$.  So any transformation in $A \in SO_3(\mathbb{R})$ may be understood as a planar rotation about an axis $v$.

\textbf{Example.} It turns out that elements of $SO_4(\mathbb{R})$ (in appropriate coordinates) look like the composition of a rotation in two separate two-dimensional subspaces.

\end{document}